\documentclass[10pt]{article}

\usepackage[utf8]{inputenc}

\usepackage[T1]{fontenc}

\usepackage{amsmath}

\usepackage{amsfonts}

\usepackage{amssymb}

\usepackage[version=4]{mhchem}

\usepackage{stmaryrd}

\usepackage{graphicx}

\usepackage[export]{adjustbox}

\graphicspath{ {./images/} }



\begin{document}

гладкой разомкнутой линией параболического вырождения $A B$, принимающего заданные значөния на әллиптической части $\sigma$ границы $\partial \Omega$ области $\Omega$ задания уравнения и на одной из двух характеристик $A C$ или $B C$, образующих гиперболическую часть $\sigma_{1}=A C \bigcup B C$ границы $\partial \Omega=\sigma \bigcup \sigma_{1}$ (см. Слешанноео $\mathrm{mu}$ па уравнение).



Т. 3. для Трикоми уравнения



\[

T u \equiv y u_{x x}+u_{y y}=0

\]



в области $\Omega$, ограниченной гладкой кривой $\sigma \subset\{(x, y)$ : $: y>0\}$ с концами в точках $A(0,0), B(1,0)$ и характеристиками $A C$ и $B C$ :



\[

A C: x=\frac{2}{3}(-y)^{3 / 2}, B C: 1-x=\frac{2}{3}(-y)^{3 / 2}

\]



впервые была поставлена и изучена Ф. Трикоми (F . Tricomi, [1], [2]).



При определенных огранитениях на гладкость заданных функций и характер поведения производной $u_{y}$ искомого решения $u$ в точках $A$ и $B$ Т. з.:



\[

\left.u\right|_{\sigma}=\varphi,\left.u\right|_{A C}=\psi

\]



для уравнения (1) редуцируется к отысканию регулярного в области $\Omega+=\Omega \bigcap\{x, y): y>0\}$ решения $u=u(x, y)$ уравнения (1), удовлетворяющего краевым условиям



\[

\begin{gathered}

\left.u\right|_{\sigma}=\varphi, \\

u_{y}(x, 0)=\alpha D_{0 x}^{2 / 3} u(x, 0)+\psi_{1}(x), 0 \leqslant x \leqslant 1,

\end{gathered}

\]



где $\alpha=$ const $>0, \psi_{1}(x)$ - однозиачио огределяется через $\psi, D_{0 x}^{2 / 3}$ - оператор дробного (в смысле Римана Лиувилля) диффференцирования порядка $2 / 3$ :



\[

D_{0 x}^{2 / 3} \tau(x)=\frac{1}{\Gamma(1 / 3)} \frac{d}{d x} \int_{0}^{x} \frac{\tau(t) d t}{(x-t)^{2 / 3}}

\]



$\Gamma(z)$ - гамма-функция Эйлера.



Решение задачи (1), (3) в свою очередь сводится $\kappa$ нахождению функции $v(x)=u_{y}(x, 0)$ из нек-рого сингулярного интегрального уравнения. Это уравнение в случае, когда $\sigma$ совнадает с нормальным контуром



\[

\sigma_{0}=\left\{(x, y):\left(x-\frac{1}{2}\right)^{2}+\frac{4}{9} y^{3}=\frac{1}{4}, y \geqslant 0\right\}

\]



имеет вид



\[

\begin{aligned}

v(x)+\frac{1}{\pi \sqrt{3}} & \int_{0}^{1}\left(\frac{t}{x}\right)^{2 / 3}\left(\frac{1}{t-x}-\frac{1}{t+x-2 x}\right) \times . \\

& \times v(t) d t=f(x),

\end{aligned}

\]



где $f(x)$ явно выражкается через $\varphi$ и $\psi$, а интеграл понимается в смысле главного значения по Коши (см. [1]-[4]).



При доказательстве единственности и существования решения T. з. наряду $\mathrm{c}$ принципом әкстремума Бицадзе (см. Смешанного типа уравнение) и методом интегральных уравнений широко используется так наз. м е т о Д $a \quad b \quad c$, сущность к-рого заключается в том, тто для данного дифференциального оператора 2 -го порядка $L$ (напр., $T$ ) с областью определения $D(L)$ строится дифференциальный оператор 1-го порядка



\[

l=a(x, y) \frac{\partial}{\partial x}+b(x, y) \frac{\partial}{\partial y}+c(x, y),(x, y) \in \Omega

\]



обладающий тем своӥством, что



\[

\int_{\Omega} l u \cdot L u d x d y \supseteqq c\|u\|^{2} \quad \forall u \in D(L)

\]



где $C=\mathrm{const}>0 ;\|\cdot\|$ - нек-рая норма (см. [5]).



Т. 3. получила обобщения как на случай уравнений смешанного типа с несколькими линиями параболиче- ского вырождения (см. [6]) так и на случай уравнений смешанного гиперболо-параболического типа (см. [7]).



Лит.: [1] Т р и к о м и Ф., О линейных уравнениях в частных производных второго порядка смешанного типа, пер.

с итал., М. - Л., 1947, [2] е г 0 ж е, Лекции по уравнениям в частных производных, плер. с италь., М., 1957; [3] Б и ц а дз е А. В., К проблеме уравнений смешанного типа, М., 1953 ; [4] е г о ж. е, Уравнения смешанного типа, М., 1959; 5] ковой газовой динамики, пер. с англ., М., 1961; [6] H a ху ші е в А. М., “Докл. АН СССР», 1966, т. 170, с. $38-40 ;$ [7] Д. ж у р а е в Т. Д., Краевые задачи для уравне

го и смешанно-составного типов, Ташкент, 1979.



A. M. Haxymes. ТРИКОМИ У РАВНЕНИЕ - дифференциальное уравнение вида



\[

y u_{x x}+u_{y y}=0

\]



к-рое является простоӥ моделью смешанного эллиптико-гиперболического типа уравнений с частными производными 2-го порядка с двумя независимыми переменными $x$, $y$ и с одной разомкнутой нехарактеристической линией параболического вырождения. Т. у. эллиптично при $y>0$, гиперболично при $y<0$ и вырождается в параболического тиша уравнение на прямой $y=0$ (см. [1]). Т. у. является прототипом у р а внения Ч а П Л Ы г ин а



\[

k(y) u_{x x}+u_{y y}=0,

\]



где $u=u(x, y)$ - функция тока плоскопараллельных устанавливающихся газовых течений, $k(y)$ и $y-$ функции скорости теqения, к-рые положительны при дозвуковой и отрицательны при сверхзвуковой скорости; $x$ - угол наклона вектора скорости (см. [2], [3]).



К краевым задачам для Т. у. сводятся многие важные проблемы механики сплошных сред, в частности, смешанные течения с образованием локальных доввуковых зон (см. [3], [4]).



Лum.: [1] Т р и к о м и Ф., О линейных уравнениях в частных производных второго порядка смешанного типа, пер. с



\begin{center}

\includegraphics[max width=\textwidth]{2023_07_09_c2a219a6daf1f48cf143g-1(1)}

\end{center}



\includegraphics[max width=\textwidth, center]{2023_07_09_c2a219a6daf1f48cf143g-1}

классы уравнений в частных производных, М., 1981.



ТриортогонА ТЕЙ - множество поверхностей в трехмерном пространстве, распадающихся на три однопараметрич. семейства таким образом, что любые две поверхности различных семейств образуют прямой угол в каждой точке их пересечения. Предполагается, что входящие в Т. с. п. поверхности являІтся регулярными, в әтом случае кривые, по к-рым пересекаются входящие в нее поверхности, являітся линиями кривизны әтих поверхностей (т е о р е м а Д ю п е н а).



Т. с. п. образуют системы координатных поверхностей в ортогональной координации пространства. Так, в сферич. системе координат Т. с. п. образуют: одно семейство сфер с общим центром в начале координат, второе семейство конусов вращения с верпиной в начале координат и с осыо, через к-рую проходят плоскости третьего семейства координатных поверхностей. С каждой Т. с. п. может быть связана нек-рая ортогональная координация пространства. Линейный элемент пространства в ортогональных координатах $u, v, w$ имеет вид



\[

d s^{2}=H_{1}^{2} d u^{2}+H_{2}^{2} d v^{2}+H_{3}^{2} d w^{2}

\]



где $H_{i}(u, v, w), i=1,2,3,-$ т. н. ф у н к ци и $Л$ а м е, для к-рых риманов тензор этой пространственной формы тождественно равен нулю. Этими фунгциями определяется Т. с. І. с точностью до движения (или отражения). С каждой регуля рной поверхностыло может быть связана Т. с. П., в состав к-рой она входит. Если задано одногараметрич. семейство регулярных поверхностей, входящее в состав Т. с. п., и если в этом семействе содержится хотя бы одна поверхность, от-





\end{document}